% !PRINTED 154
\weberpage{154}

Squaring, one obtains for the discriminant of the cubic equation (1), by \S.\,35, (6)
\weq{6}{D = -(27 b^2 + 4 a^3).}
From this one obtains the discriminant of the general function
\weq{7}{a_0 x^3 + a_1 x^2 + a_2 x + a_3,}
if for $a$ and $b$ one substitutes the expressions \S.\,34, (4), in which one replaces $a_1$, $a_2$, $a_3$ by
\[
\frac{a_1}{a_0}, \quad \frac{a_2}{a_0}, \quad \frac{a_3}{a_0}
\]
and multiplies [\S.\,46, (3)] by $a_0^4$. One thus obtains
\[
\begin{aligned}
a &= \frac{3 a_0 a_2 - a_1^2}{3 a_0^2},\\
b &= \frac{2 a_1^3 - 9 a_0 a_1 a_2 + 27 a_0^2 a_3}{27 a_0^3},
\end{aligned}
\]
and consequently
\weq{8}{D = -\frac{1}{27 a_0^2}\left\{(2 a_1^3 - 9 a_0 a_1 a_2 + 27 a_0^2 a_3)^2 + 4(3 a_0 a_2 - a_1^2)^3\right\}.}

If one works out the square and the cube, one comes back to formula (10) of the preceding paragraph; it is remarkable, however, that formula (8) apparently contains the denominator $27 a_0^2$, which cancels out in the computation.

We can proceed in the same way with the biquadratic equations. If $\alpha$, $\beta$, $\gamma$, $\delta$ are the roots of the biquadratic equation
\weq{9}{x^4 + a x^2 + b x + c = 0}
we obtain, by \S.\,37, (5)
\weq{10}{\begin{aligned}
\alpha - \beta &= v + w, & \gamma - \delta &= v - w,\\
\alpha - \gamma &= w + u, & \delta - \beta &= w - u,\\
\alpha - \delta &= u + v, & \beta - \gamma &= u - v.
\end{aligned}}
By this the discriminant of equation (9) becomes
\[
D = (v^2 - w^2)^2 (w^2 - u^2)^2 (u^2 - v^2)^2,
\]
that is, $D$ is equal to the discriminant of the cubic resolvent \S.\,37, (4)
\weq{11}{z^3 + 2 a z^2 + (a^2 - 4c) z - b^2 = 0,}
whose roots are $u^2$, $v^2$, $w^2$.

But if one forms this discriminant according to formula (10), \S.\,46, by
%CONT>>
